% 325_Hawking_FFGFT_En_ch.tex
% Doc. 325, FFGFT corpus, August 2026 — chapter file
% Included by Sources/wr_standalone_A4/325_Hawking_FFGFT_En.tex

% --- Document-local macros (collision-safe) ---
\providecommand{\K}{}\renewcommand{\K}{\textbf{[K]}}
\providecommand{\B}{}\renewcommand{\B}{\textbf{[B]}}
\renewcommand{\S}{\textbf{[S]}}
\providecommand{\Q}{}\renewcommand{\Q}{\textbf{[Q]}}
\providecommand{\X}{}\renewcommand{\X}{\textbf{[X]}}
\providecommand{\SETZ}{}\renewcommand{\SETZ}{\textbf{[POSTULATE]}}
\providecommand{\lP}{}\renewcommand{\lP}{\ell_{\mathrm{P}}}
\providecommand{\mP}{}\renewcommand{\mP}{m_{\mathrm{P}}}
\providecommand{\TP}{}\renewcommand{\TP}{T_{\mathrm{P}}}
\renewcommand{\TH}{T_{\mathrm{H}}}
\providecommand{\rs}{}\renewcommand{\rs}{r_{\mathrm{s}}}
\providecommand{\SBH}{}\renewcommand{\SBH}{S_{\mathrm{BH}}}
\providecommand{\Tmass}{}\renewcommand{\Tmass}{\widetilde{T}}
\providecommand{\Ebit}{}\renewcommand{\Ebit}{E_{\mathrm{bit}}}
% --- Colour boxes ---
\definecolor{ffgftblue}{RGB}{30,90,160}
\definecolor{warn}{RGB}{100,0,120}
\tcbset{
  base/.style={breakable,enhanced,boxrule=0.6pt,arc=3pt,
    fonttitle=\bfseries\small,coltitle=white,top=4pt,bottom=4pt,left=6pt,right=6pt},
  ffgftbox/.style={base,colback=blue!5,colframe=ffgftblue,colbacktitle=ffgftblue},
  warnbox/.style={base,colback=violet!6,colframe=warn,colbacktitle=warn}
}

\section*{Abstract}

Doc.\ 313 (Ch.\ G) established the thermodynamic frame of Hawking
radiation in FFGFT: a single KMS rule tempers Unruh, Gibbons--Hawking
and black holes jointly \K, the membrane picture receives its
thermometer \S, and the family ladder reaches the cosmos at
$\rs = R_H/2$ \S. What remained open was the \emph{microscopic
emission mechanism}: which modes escape, what the quantum carries,
why massive quanta fall back. This document works it out completely
from FFGFT-native building blocks:

\begin{enumerate}
  \item \textbf{Temperature}: The KMS rule $T = \hbar/(k_B\tau)$ with
    membrane period $\tau = 8\pi GM/c^3$ follows from $\Tmass\cdot m = 1$
    (clock compression at the membrane, Doc.\ 313 G.2). \K
  \item \textbf{Quantum}: The Hawking quantum is the bit of scale $4\pi\rs$:
    $k_B\TH = \hbar c/(4\pi\rs)$ exactly — the system-dependent bit energy
    $\Ebit = \hbar c/L$ (Doc.\ 257) evaluated at the horizon scale.
    No universal bit value is needed anywhere. \K
  \item \textbf{Selection}: Only triality-0 modes ($T_R = 0$) can be
    emitted. Confinement (Doc.\ 321) and Hawking selection are
    \emph{one} principle; absorption of the partner quantum is the
    $\mathbb{Z}_3$ projector orthogonality $P_jP_k = 0$. \B
  \item \textbf{Carrier}: The massless carrier is the $n_4 = 0$ torus mode
    (spatial winding without mass-circle winding). Massive modes
    ($n_4 \neq 0$) have finite $\Tmass$ and fall back. \K
  \item \textbf{Information}: The emitted quantum carries the sector pair
    $(k, -k)$ as an orthogonally readable quantum number ($\log_2 3$ bits) —
    plus the thermal 1-nat area bookkeeping. Fine-grained entropy remains
    constant (unitarity, Doc.\ 322). No information paradox. \K
\end{enumerate}

Two verification scripts, all assertions passed:
\texttt{325\_hawking\_ffgft\_mechanismus.py} (thermodynamic chain),
\texttt{325\_hawking\_z3\_selektion.py} ($\mathbb{Z}_3$ algebra).

% ============================================================
\section{Starting point: what Doc.\ 313 left open}

Doc.\ 313 Ch.\ G established three results: (a) the single KMS rule
$T = \hbar/(k_B\tau)$ for Unruh, Gibbons--Hawking and black holes \K;
(b) the membrane thermometer from $\Tmass\cdot m = 1$ \S; (c) the
family ladder up to the cosmos ($\rs = R_H/2$ exactly) \S. Likewise
the limit: evaporation is 56 (stellar) resp.\ 83 (supermassive) orders
of magnitude too slow for the time cycle — Hawking sharpens the
coarse-grained entropy question D(ii), it does not solve it \K.

What was missing was the \emph{microstructure} of emission:
\begin{itemize}
  \item Which modes of the $T^4/\mathbb{Z}_3$ spectrum can leave the
    horizon — and which principle selects them?
  \item What is the emitted quantum, and how is its energy related to
    the information unit of the framework (Doc.\ 257)?
  \item Why do massive quanta fall back — as a consequence of the
    geometry, not as an extra assumption?
  \item How is information conservation realized microscopically?
\end{itemize}
These four questions are answered here — each from an already
established FFGFT building block: $\Tmass\cdot m = 1$ (Docs.\ 077, 180),
system-dependent bit energy (Docs.\ 257, 302, A271--A273),
$\mathbb{Z}_3$ triality (Doc.\ 321), unitarity of the spectral operator
(Doc.\ 322).

% ============================================================
\section{Building block 1: the KMS temperature from \texorpdfstring{$\Tmass\cdot m = 1$}{T*m=1}}

Wherever the local time structure is periodic with period $\tau$,
an observer reads off temperature:
\begin{equation}
  T = \frac{\hbar}{k_B\,\tau}
  \label{eq:kms}
\end{equation}

\begin{center}
\small
\begin{tabular}{@{}lll@{}}
  \toprule
  System & Period $\tau$ & $T = \hbar/(k_B\tau)$ \\
  \midrule
  Cosmos (compact time) & $2\pi/H_0 = \num{2.9e18}$ s & $\num{2.66e-30}$ K \\
  Black hole ($M_\odot$) & $8\pi GM/c^3 = \num{1.24e-4}$ s & $\num{6.17e-8}$ K \\
  Unruh ($a = g$) & $2\pi c/a = \num{1.9e8}$ s & $\num{3.98e-20}$ K \\
  \bottomrule
\end{tabular}
\end{center}

The membrane period $\tau = 8\pi GM/c^3$ is not an extra assumption
but follows from time-mass duality: mass compresses the time structure
($\Tmass\cdot m = 1$ — the closer to the membrane, the slower the
clocks), and at the membrane the local time structure becomes periodic
with exactly this period (Doc.\ 313 G.2). \K

The temperature is thus not the result of a Bogoliubov computation on
a curved background, but a read-off rule of the clock map — the same
rule for all three systems.

% ============================================================
\section{Building block 2: the Hawking quantum as a system-dependent bit}

The central new result of this document. The system-dependent bit
energy of FFGFT (Doc.\ 257) reads $\Ebit(L) = \hbar c/L$.
Evaluated at the horizon scale:

\begin{tcolorbox}[ffgftbox, title={Core result \K}]
  \begin{equation}
    k_B \TH \;=\; \frac{\hbar c}{4\pi \rs}
    \;=\; \Ebit\!\left(L = 4\pi\rs\right)
  \end{equation}
  The thermal Hawking quantum \emph{is} the bit of scale
  $L = 4\pi\rs$ (twice the horizon circumference).
  Numerically exact (script 1): $\hbar c/(k_B\TH) = 4\pi\rs$ to
  machine precision.
\end{tcolorbox}

The system dependence of bit values (Docs.\ 257, 302, A271--A273)
is therefore no obstacle to Hawking physics but \emph{its explanation}:
the Hawking temperature is the bit temperature of the horizon scale.
A universal bit value — say, a fixed bit mass from a Landauer
evaluation at the Planck temperature — is needed nowhere; every system
carries its own bit energy, and that of the black hole at its horizon
is exactly $k_B\TH$.

\subsection{Area bookkeeping}

Emission of a quantum of energy $k_B\TH$ carries, by Clausius,
exactly $\dd S = E/(T\cdot k_B) = 1$ nat, hence
\begin{equation}
  \dd A = -4\lP^2 \quad\text{per nat}
\end{equation}
(script 1). The area loss per emission thus follows from
Bekenstein $+$ Clausius alone — no further geometric assumptions. \B

% ============================================================
\section{Building block 3: \texorpdfstring{$\mathbb{Z}_3$}{Z3} triality selection}

The algebraic core of the mechanism. The horizon carries the
$\mathbb{Z}_3$ orbifold structure of $T^4/\mathbb{Z}_3$ (Docs.\ 321, 322).
The three $\mathbb{Z}_3$ projectors
\begin{equation}
  P_k = \frac{1}{3}\sum_{j=0}^{2}\omega^{-jk}\,\tau^j,
  \qquad \omega = e^{2\pi i/3}
\end{equation}
are complete, idempotent and orthogonal (script 2):
\begin{equation}
  P_0 + P_1 + P_2 = \mathbb{1},\qquad
  P_j P_k = \delta_{jk}P_j.
\end{equation}

\begin{tcolorbox}[ffgftbox, title={Absorption \B}]
  The infalling partner quantum erases the horizon voxel exactly when
  it lies in the orthogonal $\mathbb{Z}_3$ sector:
  \begin{equation}
    P_j\,\bigl(P_k\,\psi\bigr) = 0 \quad (j \neq k)
  \end{equation}
  — exact, no tunnelling (script 2: $\|P_2(P_1\psi)\| < 10^{-15}$).
  Pair creation at the membrane separates a $(k,-k)$ sector pair;
  the inward-falling partner is absorbed by projector orthogonality,
  the outward-running one escapes provided it satisfies the triality
  condition.
\end{tcolorbox}

\subsection{Who escapes: \texorpdfstring{$T_R = 0$}{TR=0}}

The triality of a mode is its $\mathbb{Z}_3$ charge $k$.
Emittable (asymptotically free) are only combinations with
$T_R = \sum k_i \equiv 0 \pmod 3$:

\begin{center}
\small
\begin{tabular}{@{}lcc@{}}
  \toprule
  Combination & $T_R$ & emittable \\
  \midrule
  Single mode $k=0$ (photon analogue) & 0 & yes \\
  Single mode $k=1$ (quark) & 1 & no \\
  Bilinear $(1,2)$ (meson) & 0 & yes \\
  Bilinear $(1,1)$ & 2 & no \\
  Trilinear $(1,1,1)$ (baryon) & 0 & yes \\
  \bottomrule
\end{tabular}
\end{center}

\begin{tcolorbox}[ffgftbox, title={Confinement $=$ Hawking selection \B}]
  Doc.\ 321 identified confinement as triality selection $T_R = 0$.
  The same principle acts at the horizon: only colourless objects
  escape. Confinement and Hawking selection are \emph{one} principle
  of the $\mathbb{Z}_3$ structure — not two.
  The $\mathbb{Z}_3$ charge is exactly conserved in every emission
  (script 2).
\end{tcolorbox}

% ============================================================
\section{Building block 4: the massless carrier}

On $T^4 = T^3_{\mathrm{space}}\times S^1_m$ the mass of a mode is
carried by the winding $n_4$ on the mass circle ($\Tmass\cdot m = 1$).
The emittable carrier is the mode with
\begin{equation}
  n_4 = 0:\qquad \text{spatial winding without mass winding.}
\end{equation}
A massive quantum ($n_4 \neq 0$) has finite $\Tmass$; in the
compressed clock map at the membrane (Doc.\ 313 G.2) it falls back —
as a consequence of the geometry, not as an extra assumption.
For $|n_i| \leq 1$ there are $3^3 - 1 = 26$ massless neighbouring
modes against $54$ massive ones (script 2). \K

The Hawking spectrum thus consists of $n_4 = 0$ modes with $T_R = 0$:
massless, colourless, thermal at $\TH$ — the photon-like degrees of
freedom of the torus.

% ============================================================
\section{Building block 5: information conservation}

\subsection{What the quantum carries}

The emitted $T_R = 0$ quantum carries the \emph{sector pair} $(k, -k)$
as an internal quantum number. The three possible pair states are
exactly orthogonal (Gram matrix $= \mathbb{1}_3$, script 2) — which
sector was withdrawn is readable from the quantum:
$\log_2 3 \approx 1.585$ bits of sector information per quantum,
in addition to the thermal 1-nat area bookkeeping.

\subsection{Fine-grained constancy}

Doc.\ 313 G.5 phrased it as \enquote{two languages, one statement}:
locally, the radiation is correlated with the fractal membrane
structure ($S_{\mathrm{BH}}(M_\odot) = \num{1.05e77}\,k_B$ of
correlations — nothing is lost); globally, fine-grained entropy is
constant under unitary evolution (unitarity of the spectral operator,
Doc.\ 322 \K). No information paradox: the membrane is a clock-map
object, not a causal abyss.

% ============================================================
\section{The FFGFT correction and its limits}

Doc.\ 313 gives the fractal correction of the Hawking power:
\begin{equation}
  P = P_{\mathrm{std}}\,\bigl(1 - \xi\ln(M/\mP)\bigr),
\end{equation}
numerically 0.60\,\% (PBH $10^{12}$ kg) to 1.44\,\% (SMBH $10^9\,M_\odot$)
— script 1. It does not change the evaporation findings:
only primordial holes below
\begin{equation}
  M_* = \left(\frac{\tau_c\,\hbar\,c^4}{5120\,\pi\,G^2}\right)^{1/3}
      = \num{3.3e11}\ \mathrm{kg}
\end{equation}
close via Hawking within one time cycle $\tau_c = 91.9$ Gyr.
The Sun is 56, an SMBH 83 orders of magnitude too slow
(script 1) — Hawking sharpens the coarse-grained entropy question
D(ii), it does not solve it (Doc.\ 313 G.4). \K

The family ladder runs seamlessly up to the cosmos:
$\TH = T_{\mathrm{GH}}$ requires $M = c^3/(4GH_0) = \num{4.66e52}$ kg
with $\rs = R_H/2$ exactly (script 1). \K

\begin{tcolorbox}[warnbox, title={Placement of the bit result}]
  $k_B\TH = \Ebit(4\pi\rs)$ shows: Hawking physics can be formulated
  entirely with the system-dependent bit energy. Any construction that
  instead postulates a universal bit value (say, a fixed bit mass from
  $k_B T\ln 2$ at a privileged temperature) adds a parametrization
  that the temperature physics does not require — consistent with
  A271--A273: Landauer's bound is a statement about carrier
  operations, not about abstract information.
\end{tcolorbox}

% ============================================================
\section{Status overview}

\begin{center}
\begin{tabular}{@{}lc@{}}
  \toprule
  \textbf{Element} & \textbf{Status} \\
  \midrule
  KMS rule tempers Unruh/GH/BH jointly & \K\ (Doc.\ 313) \\
  Membrane period from $\Tmass\cdot m = 1$ & \K \\
  $k_B\TH = \Ebit(4\pi\rs)$ — quantum $=$ horizon bit & \K\ (new) \\
  Area quantum $-4\lP^2$/nat from Bekenstein$+$Clausius & \B\ (new) \\
  $\mathbb{Z}_3$ absorption $P_jP_k = 0$ & \B\ (new) \\
  $T_R=0$ selection; confinement $=$ Hawking selection & \B\ (new) \\
  Carrier $=$ $n_4$-free torus mode & \K\ (new) \\
  Sector info $\log_2 3$ bits orthogonally readable & \B\ (new) \\
  Information conservation (membrane $+$ unitarity) & \K\ (313/322) \\
  FFGFT power correction 0.6--1.4\,\% & \K\ (313) \\
  Greybody factors, mode-spectrum detail & \S\ (open) \\
  Back-reaction on orbifold fixed points & \S\ (open) \\
  \bottomrule
\end{tabular}
\end{center}

The Hawking mechanism in FFGFT thus stands at \K/\B; open remain only
the greybody factors (standard QFT computation on the FFGFT background)
and the back-reaction of the emission on the orbifold fixed-point
structure.

\section*{Verification scripts}

\begin{itemize}
  \item \texttt{python/Dok325\_Skripte/325\_hawking\_ffgft\_mechanismus.py} —
    thermodynamic chain: KMS, membrane thermometer, horizon bit,
    area bookkeeping, correction, evaporation, family ladder,
    information side.
  \item \texttt{python/Dok325\_Skripte/325\_hawking\_z3\_selektion.py} —
    $\mathbb{Z}_3$ algebra: projectors, absorption, triality selection,
    charge conservation, sector information, masslessness.
\end{itemize}

\section*{References}

\begin{enumerate}
  \item J.\ Pascher, Doc.\ 313 \textit{No Beginning}, Ch.\ G (Hawking on the time cycle), FFGFT corpus, 2026.
  \item J.\ Pascher, Doc.\ 321 \textit{SU(3)$_c$ from $\mathbb{Z}_3$ triality}, FFGFT corpus, 2026.
  \item J.\ Pascher, Doc.\ 322 \textit{Spectral theory and Hilbert-space embedding}, FFGFT corpus, 2026.
  \item J.\ Pascher, Docs.\ 257, 302, A271--A273 (bit energy, Landauer), FFGFT corpus.
  \item J.\ Pascher, Doc.\ 180 \textit{Derivation of $L_0$}, FFGFT corpus, 2026.
  \item S.\ W.\ Hawking, Commun.\ Math.\ Phys.\ \textbf{43}, 199, 1975.
  \item J.\ Bekenstein, Phys.\ Rev.\ D \textbf{7}, 2333, 1973.
  \item W.\ G.\ Unruh, Phys.\ Rev.\ D \textbf{14}, 870, 1976.
\end{enumerate}

\vspace{6pt}\noindent\small
\textcopyright\ 2026 Johann Pascher · \href{https://creativecommons.org/licenses/by/4.0/}{CC BY 4.0}
